\documentclass[10pt,a4paper]{article}
\usepackage[utf8]{inputenc}
\usepackage[T1]{fontenc}
\usepackage[spanish,es-noshorthands]{babel}
\usepackage[margin=3cm]{geometry}
\usepackage{amsmath,amssymb}
\usepackage{booktabs,tabularx,array}
\usepackage{enumitem}
\usepackage[most]{tcolorbox}
\usepackage{xcolor}
\usepackage{fancyhdr}
\usepackage{pdflscape}
\usepackage{graphicx}
\usepackage{caption}
\usepackage{subcaption}
\usepackage{float}
\usepackage[colorlinks=true,linkcolor=blue!50!black,urlcolor=blue!50!black]{hyperref}

% Permite cortar las líneas después de los guiones bajos (nombres de funciones y ficheros largos)
\let\guionbajo\_
\renewcommand{\_}{\guionbajo\allowbreak}
\emergencystretch=2em

\begin{document}
\sloppy

\begin{titlepage}
    \centering
    \vfill
    {\huge \textbf{PROCESAMIENTO DEL LENGUAJE NATURAL} \par}
    \vspace{2cm}
    {\Large Práctica 3: Modelos neuronales para la minería de opinión \par}
    \vspace{1cm}
    {\Large Universidad Autónoma de Madrid \par}
    \vspace{1cm}
    {\large Iván Domínguez Hernández |  ivan.dominguezh@estudiante.uam.es \par}
    \vspace{.5cm}
    {\large Lucas Miranda López | lucas.mirandal@estudiante.uam.es \par}
    \vspace{1cm}
    {\large Grupo 7461 -- Equipo 07 \par}
    \vspace{1cm}
\end{titlepage}

\tableofcontents

\newpage

\end{document}
